\chapter{Architecture, interfaces and repository reuse}\label{ch:architecture}
\section{One scientific core, several applications}
The architecture separates data access, study semantics, estimation and use.
The same immutable study definition should drive a notebook, a scheduled
service and the standalone workbench. Estimator adapters return typed
results; they do not directly approve customer actions. A result can be
computationally complete while scientifically unsupported.

\begin{figure}[htbp]\centering
\begin{tikzpicture}[node distance=8mm and 7mm,
 box/.style={draw=ink,rounded corners,fill=light,text width=42mm,align=center,minimum height=13mm,font=\small},
 flow/.style={-{Latex},thick,color=teal}]
\node[box] (data) {Bank systems\\catalog and event history};
\node[box,right=of data] (sem) {Temporal and entity layer\\snapshots, masks, lineage};
\node[box,right=of sem] (study) {Study compiler\\estimand, assumptions, folds};
\node[box,below=of study] (method) {Method adapters\\effects, panels, policies, graphs};
\node[box,left=of method] (evidence) {Evidence store\\scores, diagnostics, versions};
\node[box,left=of evidence] (apps) {Workbench and sibling apps\\decision reports and APIs};
\draw[flow](data)--(sem);\draw[flow](sem)--(study);
\draw[flow](study)--(method);\draw[flow](method)--(evidence);
\draw[flow](evidence)--(apps);
\end{tikzpicture}
\caption{Proposed architecture. Estimands, assumptions and lineage accompany results.}
\label{fig:architecture}
\end{figure}

The compiler checks required fields, intervention/population compatibility,
temporal availability, method support, dependence and action overlap.
Mechanical checks coexist with documented scientific judgments.
The execution layer schedules fold fits, scores, covariance aggregation
and diagnostics. The evidence store preserves data hashes, code/dependency
versions, seeds, configuration, assumption records and reports.
Large data and models remain in approved storage, referenced by authorized
metadata.

\section{Versioned public objects}
The supplied version 0.1.0 contracts define \emph{Study},
\emph{TemporalEvent}, \emph{EffectResult} and \emph{Graph}.
These are usable JSON Schemas with synthetic examples and selected semantic
checks. A separate version 0.2.0 research package implements the two narrow
components in Chapter~\ref{ch:construction}; the full SDK remains proposed.

Study records decision and assignment units, population/snapshot,
horizon, treatment/comparator, outcome, direction, currency, identification
route, assumptions and inference design. EffectResult preserves the
study hash, scientific status, evidence class, estimate or null,
interval semantics, independent-unit count, limitations and artifact
references. Confidence intervals, credible intervals and identified sets
must not share an unlabeled two-number field.

TemporalEvent distinguishes occurrence, recording, availability and
decision times. Graph preserves DAG/CPDAG/PAG type, lag and endpoint marks.
Checks reject late information, future historical measurements,
missing-as-zero records, unsupported point estimates, reversed intervals,
unlabeled bounds, absent assumptions, PAG edges called DAG edges,
unknown nodes and lags outside the window.

Further objects should include \emph{TrajectoryBatch},
\emph{RecoveryEnsemble}, \emph{PolicySpec}, \emph{BeliefModel},
\emph{OutcomeComponents} and \emph{EvidenceBundle}. Their required
semantics include action masks/propensities, observation/censoring masks,
state/reward timing, uncertainty draws, policy information, cash-flow
ownership and provenance. They are not already implemented schemas.

\section{Proposed API behavior}
A proposed Python workflow is:
\begin{lstlisting}
study = Study.from_json("study.json")
cohort = data.materialize(study.population, as_of=study.snapshot)
design = identify(study, cohort, graph=study.graph)
diagnostics = assess_support(design, cohort)
result = estimate(design, cohort, estimator="doubleml_adapter")
bundle = evidence.export(study, result, diagnostics)
\end{lstlisting}
These workflow names illustrate the proposed interface; they are not callable
functions in this repository. The implemented component names and actual
usage are given in Chapter~\ref{ch:construction}. Proposed service calls
return versioned artifacts and status rather than silently replacing
previous results.

\begin{center}\small
\begin{tabular}{P{76mm}P{68mm}}
\toprule Proposed endpoint & Behavior\\\midrule
\texttt{POST /v1/studies} & Validate a specification\\
\texttt{POST /v1/studies/\{id\}/runs} & Accept an idempotency key\\
\texttt{GET /v1/runs/\{id\}} & Return stage and scientific status\\
\texttt{GET /v1/runs/\{id\}/evidence} & Return authorized evidence references\\
\bottomrule
\end{tabular}
\end{center}
Data access uses approved dataset references rather than arbitrary SQL
executed with service credentials.

Errors should identify unknown treatment versions, unavailable features,
incompatible inference, unsupported actions, absent controls, unresolved
identification or numerical failure. Retries resume deterministic
checkpoints; changed inputs create a new run. Cancellation must not
publish a partial estimate as final.

\section{Cross-repository integration}
Documentation from nine local repositories was inspected and copied with
hashes \citep{localrepos}. The following are proposed interfaces,
not completed integrations.

\begin{longtable}{P{24mm}P{50mm}P{68mm}}
\toprule Repository & Proposed exchange & Constraint to preserve\\\midrule
\endhead
FinMarketing & Campaign, pricing, uplift and policy studies & Assignment versus acceptance, group contribution, horizon\\
dataml & Masks, recovery assumptions and uncertainty draws & Conditional recovery evidence; no completed table treated as error-free truth\\
BayesFilter & Filtering, smoothing and latent estimation & Event clock and numerical diagnostics; fit separate from causality\\
cardnpv & Outcome components and policy-value contrasts & Historical minus perfect-execution value; sign ID \texttt{V\_G0\_MINUS\_V\_GH}\\
SMEwallet & Offer/receipt, wallet/flow scenarios and evidence & Conservation, hidden-wallet identification limits and synthetic status\\
MacroFinance & Regime paths, structural shocks and sensitivities & Parameter identification, normalizations and shock assumptions\\
hponas & Nuisance tuning inside training partitions & No frozen causal evaluation outcomes in tuning\\
MathDev\-MCP & Equation/proof and numerical artifacts & Mathematics does not certify substantive assumptions\\
research-assistant & Literature identity and provenance & Paper organization differs from bank-data semantics\\
\bottomrule
\caption{Proposed integration map based on documentation inspection.}
\label{tab:repos}
\end{longtable}

SMEwallet's model card limits completed evidence to a synthetic prototype
and reports no promoted learned policy. Its scenarios are useful adverse
cases, not competitor-wallet ground truth. The inspected MacroFinance
milestone note discusses covariance provenance and conditional identification;
it is an earlier document, not a current overall status assessment.

For cardnpv, an adapter must map explicitly from any internal
treatment-minus-comparator convention to its canonical historical-minus-
perfect-execution direction. Test the sign on a known two-value fixture.
A numerically accurate sign reversal is a material financial error.

Integration begins with schema validation, then a synthetic adapter,
independent reference agreement and bank qualification. Repositories can
adopt contracts independently. Version mismatch must fail visibly rather
than discard fields.

\section{Deployment and operations}
Offer deployment inside the bank's controlled cloud or premises, plus local
research mode. A control service manages metadata, authorization and jobs;
workers run near data. Sensitive rows do not belong in application logs,
source-control repositories or paper-research artifacts.

Use enterprise identity, least-privilege access, encryption, audited
artifacts, retention policies and secrets management. These are proposed
requirements, to be validated in integration and security tests.
Legal-entity and regional separation applies to caches, features, models
and exports as well as source tables.

Pin adapter environments and software bills of materials. Retain exact
versions, licenses, vulnerability-review results and seeds for qualified
releases. Upgrade through reference, adverse-case and contract tests:
changed treatment coding, covariance defaults or graph endpoints can
alter a scientific conclusion.

Monitor data freshness, missingness, support, outcome delay, policy version,
graph stability and effect drift. Alarms should trigger recomputation,
restricted use or a redesigned study. They must not silently relax
identification assumptions.

\section{Model governance and current guidance}
Federal Reserve SR26-2 distributes April 2026 interagency guidance that
supersedes SR11-7 and SR21-8. It emphasizes risk-based governance,
conceptual soundness, outcomes analysis, effective challenge and vendor
responsibility \citep{fed2026}. Its supervisory applicability is specific;
it is not a global bank rule.
The scope footnote excludes generative and agentic AI models while
covering traditional quantitative and non-generative, non-agentic AI.
The semantic-assistance layer therefore needs the bank's applicable
governance assessment rather than an assertion that the guidance directly
covers every component.

PRA SS1/23 was revised in April 2026. It covers specified UK-incorporated
firms with internal-model regulatory-capital permission, with exclusions
stated in the document \citep{pra2026}. Its principles address model
lifecycle and self-assessment. Map each legal entity and use to applicable
requirements. The product supplies evidence; it cannot confer approval.
